\section{Status of the density problem}\label{sec:status}

\subsection{What is proved}
Let $T$ be an arbitrary admissible operator.  Density of
$\NA((c_0,p),\ell_2^2)$ is equivalent to $\Rec=S_{p^*}$
(Proposition~\ref{prop:reduction}), and by Lemma~\ref{lem:martintail} it
suffices to prove $\Rec=S_{p_N^*}$ for infinitely many $N$.  A first row
$f$ belongs to $\Rec$ under each of the following sufficient conditions:
\begin{enumerate}
\item[(S1)] $f$ attains its norm (Proposition~\ref{prop:reduction}(d));
\item[(S2)] $f$ belongs to the residual set $\mathcal O$ of Hausdorff continuity
points of the fibres (residual uniform recovery theorem of \cite{PrA});
\item[(S3)] $\Cset(f)=\cl\Cert^{\rm sh}(f)$; then every mate is recovered
along every sequence $f_n\to f$ (Theorem~\ref{thm:shifted}); more generally
$\Cset(f)=\mathcal K(f)$ when $a\in c_{00}$ (Proposition~\ref{prop:intrinsic});
\item[(S4)] $I$ is finite and $f$ satisfies \textup{(BT)}: finite base support,
finitely many strict non-peaks, no degenerate peaks and \textup{(MS)}, with
an \emph{arbitrary} contact set (Corollary~\ref{cor:BTrecovered}; this
contains Theorem~\ref{thm:tame} and the first row of
Section~\ref{sec:resonance});
\item[(S5)] $f$ is a limit of points of $\Rec$ along which the support
functions $\tilde r$ are asymptotically at least those of $f$
(Theorem~\ref{thm:transitivity});
\item[(S6)] $T$ is the signature-ladder operator of
Definition~\ref{def:SLD}, $I=\{1,\dots,N\}$, and $f\in\mathcal R_0$, i.e.\
$F$ is finite and the signature sets keep room (Theorem~\ref{thm:R0}); no
condition on the blocks of $f$ is needed, and $\mathcal R_0$ contains all
norm-attaining and all base-tame first rows.
\end{enumerate}
For an individual mate $g\in\Cset(f)$, $(f,g)\in\cl\NA$ whenever $g$ lies
in $\cl\Cert^{\rm sh}(f)$; this covers finite and shifted certificates
(Theorems~\ref{thm:transport}, \ref{thm:shifted}), weighted directions with
$O(\sigma)$ box tails (Theorem~\ref{thm:weighted}), locally split mates and
split-contractive operators when $I$ is finite (Theorem~\ref{thm:locsplit},
Corollary~\ref{cor:splitcontractive}), every mate that is a certificate of
radius $\gtrsim t$ up to an $O(t)$ remainder at all small scales
(Corollary~\ref{cor:ladder}), and carrier mates with non-summable gaps
(Theorem~\ref{thm:carriers}).  When $a\in c_{00}$, also every balanced
finite certificate with $\Gamma_w\le1$ is recovered
(Theorem~\ref{thm:transfer}).  When $I$ and $F$ are finite, every mate with
two-piece data (one-sided linear decompositions, possibly through infinitely
many contacts) with $\kappa_w\le1$ is recovered along engineered
norm-attaining approximants if $\Delta d_m\ge0$ in every block
(Corollary~\ref{cor:D1}), and if the blocks with $\Delta d_m<0$ satisfy the
scrambling condition \textup{(SC)} (Theorem~\ref{thm:engineered}); no
hypothesis on the number of active blocks, on the contact set or on rates of
$T$ is needed.  At \textup{(BT)} points the exact one-sided second-order
coefficients are known (Theorem~\ref{thm:onesided}).  For $\ell_2^{d+1}$ we
only prove the joint recovery of tuples of finite certificates under the
hypothesis of Theorem~\ref{thm:joint}.

On the negative side, the intrinsic mechanisms of
Sections~\ref{sec:certificates}--\ref{sec:second} do not suffice in general:
for some admissible $T$ there is a non-attaining $f$ at which the defect
$\Cset(f)\setminus\mathcal K(f)$ spans an infinite-dimensional space
(Theorem~\ref{thm:defect}), and its explicit elements are not recovered
along the canonical truncations (Proposition~\ref{prop:nonrecovery}).  This
is a defect of the intrinsic mechanisms only: that first row satisfies
\textup{(BT)}, and every one of its mates is recovered along engineered
norm-attaining approximants (Corollary~\ref{cor:exampleRecovered}).  At
every point where $a\in c_{00}$ and $K$ is finite (in particular at every
norm-attaining point) exact resonances do not exist
(Proposition~\ref{prop:exact}).  Nothing proved here points to a
counterexample; a counterexample would produce a two-dimensional Hilbert
space without Lindenstrauss' property B, answering Question~6 of Johnson and
Wolfe \cite{JW} negatively.

\subsection{What remains}
By Lemma~\ref{lem:pair} and Corollary~\ref{cor:lsc}(c), the density problem
for $p_N$ is equivalent to the following statement, in which only the defect
matters.

\begin{problem}\label{prob:main}
Let $N\ge1$, let $T$ be admissible, and let $f\in S_{p_N^*}$.  Is it true
that for every $g\in\Def(f)$ and every $\rho<1$ there are norm-attaining
$f_n\to f$ with $\rho g\in\Li_n\Cset(f_n)$?  Equivalently, is every pair
$(f,\rho g)$ with $g\in\Cset(f)\setminus\cl\Cert^{\rm sh}(f)$ a limit of
norm-attaining operators into $\ell_2^2$?
\end{problem}

A positive answer for infinitely many $N$ gives density for Mart\'in's
norm (Lemma~\ref{lem:martintail}); a negative answer for one $N$ and one
admissible $T$ gives a two-dimensional range without property B.  By
Proposition~\ref{prop:nonrecovery}, the approximants in
Problem~\ref{prob:main} cannot in general be the canonical truncations.

\emph{For the designed operator.}  For the operator $T$ of
Definition~\ref{def:SLD} and every $N$, Problem~\ref{prob:main} is
equivalent to Lemma~Z (Problem~\ref{prob:lemmaZ} and
Theorem~\ref{thm:reductionZ}): every mate of every first row must be
approximated by mates at first rows in $\mathcal R_0$, which need not attain
their norms.  What remains is base-side engineering at signature-resonant
first rows (contacts or near-contacts swallowing signature sets, as at the
first row of Section~\ref{sec:resonance}) and at first rows with infinite
base support (Remark~\ref{rem:lemmaZ}).

\emph{For an arbitrary admissible $T$.}  The following classes are open.

\begin{problem}[Negative $d$-mismatch]\label{prob:resonant}
Let $I$ be finite, $f\in S_{p^*}$ with $F$ finite, and let $g\in\Cset(f)$
carry two-piece data with $\kappa_w\le1$ such that the blocks with
$\Delta d_m<0$ do not satisfy \textup{(SC)}.  Is $(f,g)\in\cl\NA$?
\end{problem}

Theorem~\ref{thm:engineered} settles $\Delta d\ge0$ (for any number of
active blocks and without steering) and $\Delta d<0$ under \textup{(SC)};
Remark~\ref{rem:withoutSC}(a) gives a quantitative partial answer, and
Remark~\ref{rem:withoutSC}(b) indicates that \textup{(SC)} can fail for
admissible $T$ (sketch).  The explicit switching mates of
Theorem~\ref{thm:defect} are $d$-neutral and are recovered.

\begin{problem}[Generic supports]\label{prob:generic}
Let $F$ be finite and let some $Q_m$ be infinite, or \textup{(MS)} fail, or
degenerate peaks exist.  Is every $g\in\Cset(f)$ a limit of mates
$g_n\in\Cset(f_n)$, $f_n\to f$, carrying two-piece data to which
Theorem~\ref{thm:engineered} applies?
\end{problem}

At such $f$ the lower bound of Theorem~\ref{thm:onesided} is not available
(the block expansion is multiscale), and mates need not have two-piece data;
mixed mates combining deep weighted parts with switching parts are not
covered either.

\begin{problem}[Second-order defect]\label{prob:secondorder}
Is there $f$ with $a\in c_{00}$ and a balanced finite certificate
$g\in\Cset(f)\cap S(f)$ with $\Gamma_w(g)>1$?  By
Proposition~\ref{prop:lowerbound} such an $f$ must have degenerate peaks,
infinitely many strict non-peaks, failure of (MS), or infinitely many
contacts; in the last case a finite-dimensional analogue suggests that the
second-order coefficient can drop below $\Gamma_w$ (Remark~\ref{rem:kinks}).
At \textup{(BT)} points the question does not affect recoverability:
Corollary~\ref{cor:BTrecovered} recovers every mate, and the relevant
coefficients are the one-sided $\gamma^\pm\le\Gamma_w$
(Theorem~\ref{thm:onesided}).
\end{problem}

\begin{problem}[Approximate resonances]\label{prob:approx}
Do switching mates carried, on both sides of $t=0$, by near-threshold
strict non-peaks with summable gaps or by weak peaks belong to
$\cl\Cert(f)$, or at least to the closure of the mates recovered in
Section~\ref{sec:engineered}?  There is no linear obstruction for them
(Remark~\ref{rem:summablegaps}).  One way to prove non-membership in
$\cl\Cert(f)$ is to exhibit a sequence $f_n\to f$ with
$g\notin\Li_n\Cset(f_n)$ (Corollary~\ref{cor:lsc}).
\end{problem}

Such scale-dependent switching is the main remaining mechanism for an
arbitrary admissible $T$.  We note that operators $T$ have been proposed whose
near-threshold peak carriers are one-sided with bounded conversion capacity;
for them no recovery argument is known (the available arguments leave a
capacity requirement that does not improve as $\rho\to1$), and no
counterexample either.

\begin{problem}[Infinite base support]\label{prob:infiniteF}
Extend Theorems~\ref{thm:onesided}, \ref{thm:engineered} and~\ref{thm:R0}
to first rows with $a\notin c_{00}$.  All results of
Sections~\ref{sec:engineered} and~\ref{sec:designed} use that $F$ is finite
(no flips on $F$ below a fixed scale).
\end{problem}

\begin{problem}[Universality of the defect]\label{prob:universal}
Does \emph{every} admissible $T$ (in particular every $T$ produced by the
proof of Mart\'in's Lemma~B) admit a first row with nonempty defect?  The
linear obstruction requires a first row with few strict non-peaks, which an
adversarial $T$ may prevent; for the designed operator of
Section~\ref{sec:designed} the defect is empty at every first row in
$\mathcal R_0$ (Corollary~\ref{cor:nodefectR0}).
\end{problem}

Finally we note a gap between necessary and sufficient conditions in the
reduction itself.  If $g_1,\dots,g_d\in\Cset(f)$, then
$(g_1,\dots,g_d)/\sqrt d\in\Cset_d(f)$, and density into $\ell_2^{d+1}$
implies that such tuples are recovered along a common norm-attaining
sequence (by the argument of Lemma~\ref{lem:pair}, with rotations of
$\ell_2^{d+1}$).  But we do not
know whether density into $\ell_2^2$ alone forces the existence, for each
$f$, of one norm-attaining sequence along which the whole fibre $\Cset(f)$ is
recovered (condition (a) of Corollary~\ref{cor:finite}); only the
sufficiency of that condition is used.

